\chapter*{General Introduction}
\addcontentsline{toc}{chapter}{General Introduction}

The financial relationship between large corporate actors and their commercial partners has become one of the most sensitive points in modern banking operations. Suppliers often need liquidity before invoice maturity. Buyers seek longer payment terms without weakening their commercial network. Banks, positioned between these two forces, must transform trade data into controlled financing operations while preserving security, traceability, regulatory discipline, and operational speed. This is the business terrain in which Supply Chain Finance operates.

Supply Chain Finance, commonly referred to as SCF, covers a family of financing mechanisms designed to optimize working capital across a commercial chain. Its purpose is not limited to granting credit. It restructures the timing of cash movement between anchors, counterparties, and financial institutions. In an SCF program, a bank can rely on the commercial strength of an anchor, the validity of invoices or other underlying instruments, and a set of negotiated rules to provide financing to suppliers, buyers, or distributors. Trade Finance completes this picture by addressing broader instruments and documentary operations linked to commercial exchange. Together, these domains form a demanding digital environment: the system must understand business roles, legal documentation, product definitions, fee rules, approval workflows, and transaction lifecycles.

This end-of-studies internship took place at Adria, from February 2026 to July 1st, 2026, within the Supply Chain Finance Platform Team. The project was conducted in an Agile Scrum context, where development tasks were organized through Jira tickets, source code was managed with Bitbucket, and each contribution followed a pull request review process led by the technical lead. The first two weeks were dedicated to remote integration and training, with a dual focus on the business logic of the Adria architecture and the technical stack used by the team, especially Java and Spring Boot. After this onboarding period, the work continued on-site as part of the daily development rhythm of the platform.

The internship was not limited to observation or peripheral tasks. It was situated inside an active software delivery cycle. The team was working on a Supply Chain Finance and Trade Finance platform intended for banking use cases in Morocco and North Africa. The platform digitizes the financing of trade receivables and payables between corporate anchors and their counterparties, including suppliers, distributors, and SMEs. Its functional perimeter includes reverse factoring, dynamic discounting, receivables finance, payables finance, and distributor finance. These products are configured through programs that bind together an anchor, counterparties, financing limits, fees, disbursement rules, repayment rules, and governance states.

From a technical perspective, the system follows a microservices-oriented architecture built on Adria standards. These standards provide pre-configured Spring Cloud services and reusable infrastructure patterns that support new products developed inside the company. The platform includes a core SCF service, a gateway service, a foundation data library, a properties service, a registry service, a configuration service, and supporting services for audit, storage, authentication, and document conversion. The backend stack relies mainly on Spring Boot, Spring Security, JPA/Hibernate, PostgreSQL, Eureka, Spring Cloud Config, MinIO, Gotenberg, RabbitMQ, and Keycloak. The frontend is implemented as a React application using Vite, TypeScript, React Router, React Query, Axios, and a component system based on shadcn/ui and Radix UI.

Security and governance represent a major axis of the platform. Authentication is based on Keycloak and JWT tokens, while role extraction determines whether the user operates from a bank, anchor, or counterparty view. Backend access is protected through role-based authorization. Business validation follows the maker/checker principle, a common banking governance pattern where the user who creates or modifies a record cannot approve the same operation. This principle is implemented through approval statuses, pending actions, repository specifications, and rollback mechanisms based on original data snapshots. The same concern for control appears in the gateway layer, which handles request routing, antivirus scanning for uploaded documents, brute-force protection, maintenance mode, request logging, and correlation identifiers.

Within this broader platform, the academic subject of the internship is the Disbursement Module. Its development runs in parallel with the team's sprint work because the main platform had not yet reached the sprint milestone where the module would be integrated natively into the architecture. This situation shaped the project in a very concrete way. The module had to be studied, designed, and progressively prepared while the platform itself was still evolving from Sprint 2 toward Sprint 8. Functional refinement sessions with the Product Owner played a central role in this process: the business rules were discussed, diagrams were prepared, and the proposed flows were evaluated and adjusted according to the expected target architecture.

The central problem addressed by this report is therefore twofold. The first dimension is functional: how can a banking platform organize Supply Chain Finance operations in a way that reflects the realities of anchors, counterparties, invoices, programs, fees, disbursements, repayments, and approvals? The second dimension is technical: how can such a platform be implemented with a modular architecture that remains secure, maintainable, auditable, and adaptable to several banking contexts? The Disbursement Module stands at the intersection of these two dimensions because it converts validated financing intent into an operational payment flow, while depending on upstream entities such as programs, products, invoices, fee catalogues, and eligibility rules.

This report presents the internship work through six chapters. The first chapter introduces the host organization, the business context of Supply Chain Finance and Trade Finance, the project environment, and the methodology followed during the internship. The second chapter studies the global architecture and technical stack, with attention to the microservices, gateway, security, data storage, and deployment environment. The third chapter examines the SCF domain model and backend implementation, including entities, maker/checker governance, counterparty management, product definitions, program configuration, invoice upload, and transaction inquiry. The fourth chapter focuses on frontend implementation and user flows, from authentication and role-based navigation to invoice, document, early payment, and finance disbursement screens. The fifth chapter is dedicated to the Disbursement Module as the PFE project, covering its motivation, specification, design, implementation choices, and validation. The final chapter evaluates the work performed during the internship, reflects on the acquired skills, and identifies possible future improvements.

The ambition of this document is not merely to list technologies or reproduce development tasks. It aims to analyze the construction of a banking software platform as an engineering problem shaped by business constraints, security expectations, team methodology, and architectural decisions. The work described here was carried out inside a real delivery environment, with all the friction that such an environment implies: evolving requirements, sprint priorities, review cycles, integration constraints, and the need to align academic objectives with production-oriented software practices.
